\documentclass[a4paper,11pt]{article}

\usepackage[utf8]{inputenc}
\usepackage[T1]{fontenc}
\usepackage[ngerman]{babel}
\usepackage{graphicx}
\usepackage{booktabs}
\usepackage{float}

\usepackage[scaled]{helvet}
\renewcommand{\familydefault}{\sfdefault}

\usepackage[a4paper, top=2cm, bottom=2cm, left=2cm, right=2cm]{geometry}

\usepackage{setspace}
\setstretch{1.5}

\setlength{\parindent}{0pt}
\setlength{\parskip}{6pt}

\usepackage{microtype}
\sloppy
\hyphenpenalty=1000
\tolerance=3000

\usepackage{titlesec}
\titleformat{\section}{\normalfont\fontsize{12}{14}\bfseries}{\thesection}{1em}{}
\titleformat{\subsection}{\normalfont\fontsize{12}{14}\bfseries}{\thesubsection}{1em}{}

\usepackage[
  colorlinks=true,
  linkcolor=black,
  filecolor=black,
  urlcolor=blue
]{hyperref}

\usepackage{fancyhdr}
\pagestyle{fancy}
\fancyhf{}
\renewcommand{\headrulewidth}{0pt}
\fancyfoot[C]{\thepage}

\usepackage{listings}
\usepackage{xcolor}

\definecolor{codegreen}{rgb}{0,0.6,0}
\definecolor{codegray}{rgb}{0.5,0.5,0.5}
\definecolor{codepurple}{rgb}{0.58,0,0.82}
\definecolor{backcolour}{rgb}{0.95,0.95,0.92}

\lstdefinestyle{pythonstyle}{
    backgroundcolor=\color{backcolour},
    commentstyle=\color{codegreen},
    keywordstyle=\color{magenta},
    numberstyle=\tiny\color{codegray},
    stringstyle=\color{codepurple},
    basicstyle=\ttfamily\footnotesize,
    breaklines=true,
    captionpos=b,
    keepspaces=true,
    numbers=left,
    numbersep=5pt,
    showstringspaces=false,
    tabsize=4
}
\lstset{style=pythonstyle}

\begin{document}

\begin{titlepage}
    \thispagestyle{empty}
    \centering
    \vspace*{5cm}
    {\Huge\bfseries Projekt: Data Analysis DLBDSEDA02\_D \par}
    \vspace{1cm}
    {\Large Endprodukt \par}
    \vspace{0.3cm}
    {\large Aufgabe 1: NLP-Techniken anwenden, um eine Textsammlung zu analysieren \par}
    \vspace{0.5cm}
    {\large Studiengang: Angewandte Künstliche Intelligenz \par}
    \vspace{0.5cm}
    {\large Sven Behrens \par}
    \vspace{0.5cm}
    {\large Matrikelnummer: 42303511 \par}
    \vspace{0.5cm}
    {\large \today \par}
    \vspace{0.5cm}
    {\large Tutor: Prof.~Dr.~Frank Passing \par}
\end{titlepage}

\pagenumbering{arabic}
\setcounter{page}{1}

\section{GitHub-Repository}

Der vollständige, versionskontrollierte Quellcode ist öffentlich verfügbar unter:

\url{https://github.com/svenb23/DataAnalysis}

\section{Technische Beschreibung}

\subsection{Voraussetzungen}
\begin{itemize}
    \item Python 3.13
    \item Git
    \item CFPB Consumer Complaint Database (CSV-Download von
          \url{https://www.consumerfinance.gov/data-research/consumer-complaints/})
\end{itemize}

\subsection{Installation}
\begin{lstlisting}[language=bash]
# Repository klonen
git clone https://github.com/svenb23/DataAnalysis.git
cd DataAnalysis

# Virtuelle Umgebung erstellen und aktivieren
python -m venv venv
venv/Scripts/activate        # Windows
# source venv/bin/activate   # macOS/Linux

# Abhaengigkeiten installieren
pip install -r requirements.txt

# spaCy-Sprachmodell laden
python -m spacy download en_core_web_sm
\end{lstlisting}

\subsection{Datenvorbereitung}
Die Datei \texttt{complaints.csv} muss von der CFPB-Website heruntergeladen und im Ordner
\texttt{data/} abgelegt werden.

\subsection{Ausführung}
Die fünf Skripte werden sequentiell ausgeführt:

\begin{lstlisting}[language=bash]
python src/01_load_data.py        # Daten laden und filtern
python src/02_preprocessing.py    # Texte vorverarbeiten
python src/03_vectorization.py    # TF-IDF und Word2Vec
python src/04_topic_modeling.py   # LDA, NMF, K-Means
python src/05_visualization.py    # Plots erzeugen
\end{lstlisting}

Die Ergebnisse (Plots, CSVs) werden im Ordner \texttt{results/} gespeichert.

\subsection{Projektstruktur}
\begin{lstlisting}
DataAnalysis/
  data/              # Roh- und verarbeitete Daten
  src/               # Python-Skripte (Pipeline)
    01_load_data.py
    02_preprocessing.py
    03_vectorization.py
    04_topic_modeling.py
    05_visualization.py
  results/           # Ergebnis-Plots und CSVs
  reports/           # LaTeX-Berichte (Phase 1-3)
  requirements.txt   # Python-Abhaengigkeiten
  README.md
\end{lstlisting}

% ============================================================
% ANHANG: QUELLCODE
% ============================================================
\newpage
\appendix
\section*{Anhang: Quellcode}

\subsection*{01\_load\_data.py}
\lstinputlisting[language=Python]{../src/01_load_data.py}

\subsection*{02\_preprocessing.py}
\lstinputlisting[language=Python]{../src/02_preprocessing.py}

\subsection*{03\_vectorization.py}
\lstinputlisting[language=Python]{../src/03_vectorization.py}

\subsection*{04\_topic\_modeling.py}
\lstinputlisting[language=Python]{../src/04_topic_modeling.py}

\subsection*{05\_visualization.py}
\lstinputlisting[language=Python]{../src/05_visualization.py}

\end{document}
